% picture: "What the system is" -- the effective picture, built model by model (models 02+09, 11+01, 16, 17),
% then the ranked observables. ASD-STE100 rules as in intro.tex. Revised 2026-10-07.
% Numbers: scratchpad fact sheets of 2026-10-07 (latest scored round of each card). Derived values are labeled "derived".
% Float calls sit one subsection before their text so that the floats land near it.
\section{What the system is}\label{sec:picture}

This section builds the effective picture, one model at a time (Fig.~\ref{fig:picture}). For each equation we give its meaning, a physical picture, its measurement from the logs and the best evidence; Sec.~\ref{sec:models} gives the full evidence. Values are for regime III, with 95\% confidence intervals (CIs) in brackets, unless we state otherwise. Table~\ref{tab:words} explains each special word in one line; the Introduction explains the call clock, the read-out call and the in-flight placebo.

% glossary: "Words we use" box for Sec. II (added 2026-10-07). One plain line per entry.
% The inline first-use explanations in the abstract, the Introduction and Sec. II are still the primary definitions.
\newcommand{\gw}[2]{\textbf{#1}: #2\par\vspace{0.6pt}}
\begin{table*}[!tbp]
\caption{Words we use, in plain terms.}
\label{tab:words}
\fcolorbox{nullc}{white}{%
\begin{minipage}[t]{0.465\textwidth}\scriptsize\raggedright
\gw{Agent}{one language model inside a scaffold; it uses a computer and a group chat.}
\gw{Scaffold}{the program that runs an agent: it builds the model's input and carries out its output.}
\gw{Model call (call)}{one turn: the model reads its context and writes one output (a message, an action or a wait).}
\gw{Context window}{the text that a call reads; the agent's short-term memory.}
\gw{Forced erasure (wipe, reset)}{the scaffold empties the context window at a fixed size (41 records in regime III); notes and files survive.}
\gw{Call clock}{time counted in an agent's own calls, not in minutes. Ten idle minutes between two calls add nothing.}
\gw{Read-out call}{the first call whose context holds a given message: the moment the agent sees it.}
\gw{In flight}{posted, but not yet in the reader's context. Our placebo.}
\gw{Named message}{a message that mentions the reader by name.}
\gw{Goal period (\#1--\#51)}{the time for one shared goal, usually a week.}
\gw{Kickoff}{the message that announces a goal.}
\gw{Room}{a separate chat channel.}
\gw{Regime I, II, III}{the three eras of the scaffold: one room in daily sessions; several rooms in sessions; continuous, own room only, forced erasures.}
\gw{Scheduler}{the operator's daily timetable that turns agents on and off.}
\gw{Natural experiment (NE\#\#)}{a dated step change in the setup (scaffold, roster, rooms, operator), tested before vs.\ after; not a goal period, and often spans two.}
\end{minipage}}\hfill
\fcolorbox{nullc}{white}{%
\begin{minipage}[t]{0.465\textwidth}\scriptsize\raggedright
\gw{Nudge}{an automatic prompt to an idle agent.}
\gw{Unit}{a goal period, or the part of one between two scaffold changes.}
\gw{Impostor}{a shared input that fakes influence between agents.}
\gw{Null}{the simplest model with no effect; a claim must beat it.}
\gw{Spin}{a part with a few states; here an agent that talks, works or waits.}
\gw{Field}{an input that acts on each agent directly: timetable, goal, room, style.}
\gw{Coupling}{the effect of one agent on another, through a message that it reads.}
\gw{Quench}{a sudden change of the field, such as a kickoff.}
\gw{Loop gain $g$}{new messages that one message causes directly. Below 1, echoes die out.}
\gw{Fano factor}{variance of a count over its mean: 1 for purely random events, more for bursts.}
\gw{Collective mode}{a pattern in which many agents move together.}
\gw{Marchenko--Pastur edge}{the largest correlation that chance alone gives in a record of this size.}
\gw{Well / kick}{an agent's slow resting point / the short push of one message read.}
\gw{Reserved data}{goal periods set aside before exploration, for confirmation tests.}
\gw{M02, \ldots, M17}{a physics model, named by its folder number (Sec.~\ref{sec:method}).}
\end{minipage}}
\end{table*}

% fig_picture: schematic of the effective picture (Sec. II). TikZ, few-word labels.
% Spin states: red/blue (Ising up/down). Fields: orange. Read-out links: dark gray.
\definecolor{spinup}{HTML}{CC3311}
\definecolor{spindn}{HTML}{0072B2}
\begin{figure*}[t]
\centering
\resizebox{\textwidth}{!}{%
\begin{tikzpicture}[font=\sffamily\small,
  agent/.style={circle,draw=inkc,line width=0.5pt,minimum size=5.5mm,inner sep=0pt},
  up/.style={agent,fill=spinup!85},
  dn/.style={agent,fill=spindn!85},
  link/.style={-{Stealth[length=1.6mm]},draw=inkc!70,line width=0.5pt},
  field/.style={-{Stealth[length=2.4mm]},draw=fieldc,line width=2.2pt},
  lab/.style={font=\sffamily\footnotesize,text=inkc}]
% ---------------- (a) the swarm
\node[lab,anchor=west,font=\sffamily\small\bfseries] at (-0.2,4.6) {(a) swarm};
\node[draw=fieldc,fill=fieldc!15,rounded corners=2pt,minimum width=6.6cm,minimum height=0.55cm,lab] (F) at (3.3,3.95) {fields: scheduler $\cdot$ kickoff $\cdot$ room $\cdot$ style};
\draw[dashed,rounded corners=6pt,draw=nullc] (0,0) rectangle (3.1,2.75);
\draw[dashed,rounded corners=6pt,draw=nullc] (3.5,0) rectangle (6.6,2.75);
\node[lab,text=nullc] at (1.55,-0.3) {room A};
\node[lab,text=nullc] at (5.05,-0.3) {room B};
\node[up] (a1) at (0.6,2.1) {}; \node[up] (a2) at (1.7,2.25) {}; \node[dn] (a3) at (2.55,1.5) {};
\node[up] (a4) at (0.75,0.75) {}; \node[up] (a5) at (1.85,0.6) {};
\node[dn] (b1) at (4.1,2.1) {}; \node[dn] (b2) at (5.2,2.25) {}; \node[up] (b3) at (6.05,1.4) {};
\node[dn] (b4) at (4.25,0.7) {}; \node[dn] (b5) at (5.35,0.6) {};
\foreach \x in {0.8,2.3,4.3,5.8} \draw[field] (\x,3.62) -- (\x,2.85);
\draw[link] (a1) -- (a2); \draw[link] (a2) -- (a3); \draw[link] (a4) -- (a5);
\draw[link] (b1) -- (b2); \draw[link] (b5) -- (b3); \draw[link] (b4) -- (b5);
\draw[link,densely dotted] (a3) -- (b4);
\node[lab,align=center] at (3.3,-0.85) {named read: weak, $g\approx0.13$};
% ---------------- (b) the call clock
\begin{scope}[xshift=7.6cm]
\node[lab,anchor=west,font=\sffamily\small\bfseries] at (-0.2,4.6) {(b) call clock};
\node[lab,anchor=east] at (0.35,3.3) {sender};
\node[lab,anchor=east] at (0.35,1.6) {reader};
\draw[->,draw=nullc] (0.45,0.6) -- (4.9,0.6) node[lab,anchor=north east] at (4.9,0.55) {time};
\fill[nullc!35] (0.55,3.1) rectangle (1.9,3.5);
\draw[-{Stealth[length=2mm]},draw=spindn,line width=1pt] (1.9,3.3) -- (1.9,2.45);
\node[lab,anchor=south] at (1.9,3.5) {post};
% reader calls
\fill[nullc!35] (0.55,1.4) rectangle (1.35,1.8);
\fill[nullc!60] (1.5,1.4) rectangle (2.75,1.8);
\fill[spindn!80] (2.9,1.4) rectangle (3.75,1.8);
\fill[nullc!35] (3.9,1.4) rectangle (4.7,1.8);
\draw[densely dotted,draw=inkc] (1.9,2.45) -- (1.9,1.85);
\draw[-{Stealth[length=2mm]},draw=spindn,line width=1pt] (1.9,2.45) to[out=-10,in=100] (3.3,1.85);
\node[lab,anchor=north,align=center] at (2.12,1.35) {in flight:\\no effect};
\node[lab,anchor=north,align=center] at (3.4,1.35) {read-out:\\acts};
\end{scope}
% ---------------- (c) the well
\begin{scope}[xshift=13.0cm]
\node[lab,anchor=west,font=\sffamily\small\bfseries] at (-0.2,4.6) {(c) one agent};
\draw[draw=inkc,line width=0.9pt,domain=-2:2,smooth,variable=\x] plot ({\x+2.2},{0.42*\x*\x+0.8});
\fill[fieldc] (2.2,0.8) circle (0.07);
\node[lab,anchor=north] at (2.2,0.72) {well centre};
\shade[ball color=spindn!70] (2.95,1.12) circle (0.17);
\draw[-{Stealth[length=2mm]},draw=spindn,line width=1.1pt] (3.12,1.3) -- (3.7,1.95);
\node[lab,anchor=south] at (3.75,2.2) {kick $\approx$7 calls};
\draw[-{Stealth[length=2mm]},draw=inkc,line width=0.8pt] (2.7,1.22) to[bend left=25] (2.35,0.98);
\node[lab,anchor=south] at (2.0,2.75) {well $\approx$100 calls};
\draw[draw=inkc!60,line width=0.4pt] (2.0,2.72) -- (2.0,1.0);
\end{scope}
% ---------------- (d) memory
\begin{scope}[xshift=18.4cm]
\node[lab,anchor=west,font=\sffamily\small\bfseries] at (-0.2,4.6) {(d) memory};
\foreach \y/\c in {0.7/35,1.1/45,1.5/55,1.9/65,2.3/75,2.7/85} \fill[spindn!\c!white] (0.3,\y) rectangle (2.5,\y+0.32);
\draw[draw=inkc,line width=0.6pt] (0.2,0.6) rectangle (2.6,3.12);
\node[lab,anchor=south] at (1.4,3.15) {context window};
\node[lab,anchor=west,align=left] at (2.75,2.55) {erase:\\$-39\%$ work};
\node[lab,anchor=west,align=left] at (2.75,1.3) {notes,\\search: $\approx0$};
\node[lab] at (1.4,0.25) {no group store};
\end{scope}
\end{tikzpicture}}
\caption{The effective picture. (a)~Agents are spins (two colors: two action states). Strong fields (orange) act on every agent; weak links (gray) act only through named messages that a call reads, mostly inside a room (dashed boxes). (b)~A message posted while the reader's call runs is in flight and has no effect; it acts at the reader's next call, the read-out call. (c)~Each agent sits in a slow well set by its fields; a read kicks it, and the kick decays about 15 times faster than the well (H130). (d)~The context window holds the working state: erasing it costs work (H15), while memory notes and history search carry almost no measured value (H87, H84). Look at the arrow widths in (a): fields, not links, set most of the state.}
\label{fig:picture}
\end{figure*}

% table_monitor: ranked observables for operators (Sec. II). Long form: appendix (tab:observables).
% Ranks set 2026-10-07 from the cards' evidence (pre-registration, strongest null, unfitted check, periods, fragile flag).
% Layout 2026-10-07 (Vivian: "hard to read"): four plain-language groups; use case split into warning sign and action;
% untested actions marked with a dagger; F and U drawn as three-dot meters. Numbers unchanged.
\newcommand{\rkdot}[2]{\ifnum#2>#1\relax\draw[nullc!80,line width=0.35pt,fill=white] (#2*4.2pt,0) circle (1.45pt);\else\fill[inkc] (#2*4.2pt,0) circle (1.6pt);\fi}
\DeclareRobustCommand{\rk}[1]{\tikz[baseline=-0.6ex]{\rkdot{#1}{1}\rkdot{#1}{2}\rkdot{#1}{3}}}
\newcommand{\grp}[1]{\multicolumn{6}{@{}l}{\rule{0pt}{2.6ex}\textsc{#1}}\tabularnewline[1pt]}
\newcommand{\nt}{$^\dagger$}
\begin{table*}[!tp]
\caption{Quantities for operators to watch, grouped by what they measure and sorted by faithfulness. Values are exploratory and for regime III unless stated. \emph{F} (faithfulness): \rk{3} a pre-registered test beats the strongest null and a statistic the model did not fit agrees, in many periods; \rk{2} it beats the strongest null, but the test is post hoc or fragile, or covers few periods; \rk{1} post hoc in one period, or not measured. \emph{U} (usefulness to an operator): \rk{3} an alarm with a measured false-alarm rate, or a lever with a measured effect; \rk{2} an early warning whose action is untested; \rk{1} a description only. \nt\,The action has not been tested. No quantity is confirmed on reserved data yet.}
\label{tab:monitor}
\footnotesize\renewcommand{\arraystretch}{1.12}
\begin{tabular}{@{}p{0.175\textwidth}p{0.215\textwidth}p{0.205\textwidth}p{0.27\textwidth}@{\hspace{5pt}}c@{\hspace{5pt}}c@{}}
\toprule
Quantity & Now & Warning sign & What to do & F & U\tabularnewline\midrule
\grp{Couplings and echoes: do agents set each other off?}
\rr \textbf{Name effect}: talk after a named read over an unnamed one & \rr $\times$19: 0.079 vs 0.004 per read (H67); content $\times$5 (H50) & \rr falls toward 1: names stop steering attention & \rr To move one agent, name it: a named nudge adds about 1 active minute per 30\,min (H30, fragile). & \rk{3} & \rk{3}\tabularnewline[2.6pt]
\rr \textbf{Loop gain} $g$: extra messages per message (Eq.~\ref{eq:fano}) & \rr 0.13 (25 units); upper bound $<0.40$ in 66/66 units (H67) & \rr rises toward 0.4: the swarm starts to echo itself & \rr Cut posts to several rooms, or split rooms, before a cascade forms.\nt & \rk{3} & \rk{2}\tabularnewline[2.6pt]
\rr \textbf{Amplification} $\chi=1/(1-g)$: total response to one push & \rr 1.15 (derived); $\le1.31$ in 33/33 periods (H51) & \rr above 2: a push to one agent reaches the group & \rr Send broad messages with care.\nt & \rk{3} & \rk{2}\tabularnewline[2.6pt]
\rr \textbf{Fano ratio} $r_F=\Phi_{\rm obs}/\Phi_{\rm pred}$ (Eq.~\ref{eq:fano}) & \rr 1.01 [0.94, 1.09], 18 units; regime I 1.34 [1.18, 1.52] (H111) & \rr leaves 0.8--1.2: a drive or coupling outside the loop acts & \rr Find it before you trust $g$. Needs only talk counts. & \rk{3} & \rk{2}\tabularnewline[2.6pt]
\rr \textbf{Idea branching} $\hat R$: new adopters per adopter & \rr median 0.22 (0.09--0.40); upper bound $<1$ in 32/32 periods (H34) & \rr upper bound reaches 1: one idea can sweep the swarm & \rr Check the idea's source.\nt & \rk{3} & \rk{2}\tabularnewline[2.6pt]
\rr \textbf{Call clock} $\eta$: 0 if calls set time, 1 if minutes do & \rr $-0.11$ [$-0.16$, $-0.06$]; regime I $+0.51$ [0.41, 0.61] (H40) & \rr above 0.5: wall time sets the coupling, as in regime I & \rr Measure the per-call quantities again before you use them. & \rk{3} & \rk{1}\tabularnewline[2.6pt]
\grp{Fields: what reaches every agent at once?}
\rr \textbf{Kickoff alignment} and topic alarm: day-1 content against the kickoff & \rr own kickoff first in 18/33 periods (H54); alarm hits 82--85\% of kickoffs, 7--9\% of placebos (H36) & \rr alarm silent 30\,min after a kickoff: the goal did not land & \rr Send the kickoff again, naming a target.\nt & \rk{3} & \rk{3}\tabularnewline[2.6pt]
\rr \textbf{Change monitor}: daily log schema, counts and topic & \rr false alarms 0.085 [0.04, 0.17] per day; hits 0.57 of kickoffs, 6/6 room changes (H74, fragile) & \rr it fires: the platform or a drive changed & \rr Read the changelog, then measure everything again. & \rk{2} & \rk{3}\tabularnewline[2.6pt]
\rr \textbf{Field share} $f$: co-activation explained by the daily start and stop & \rr 0.72 median, 0.83 after trim; regime I 0.11 (H38, fragile) & \rr below 0.5: agents act together without the schedule & \rr Find the new path of coupling, such as a room merge.\nt & \rk{2} & \rk{2}\tabularnewline[2.6pt]
\rr \textbf{Day-edge irreversibility}: share of time-asymmetry in the first and last 30\,min & \rr 0.80--1.00 of the excess in 12--25\% of the time (H76) & \rr excess appears inside the day: a new drive acts there & \rr Find the drive. & \rk{2} & \rk{1}\tabularnewline[2.6pt]
\grp{Memory: what does an agent keep, and for how long?}
\rr \textbf{Erasure cost}; value per bit $\kappa_C$ & \rr $-39\%$ commits over 10 calls (H15); $\kappa_C=5.2$ [3.8, 7.9] (H87); a 20-call cap keeps 0.88 [0.85, 0.91] of output per call (H44) & \rr a planned cut of the context cap, or a forced reset & \rr Read the cost first. Inside an idle trap (a long run of waits), a forced erasure lifts escape from 0.47 to 0.84 (H16). & \rk{2} & \rk{3}\tabularnewline[2.6pt]
\rr \textbf{Well depth} $\Delta\rho$: extra loss of old content at a kickoff & \rr $+0.26\pm0.04$ (s.e.), 16/18 kickoffs (H97) & \rr near 0 after a new goal: the old goal still holds & \rr Expect old content on day 1 (H82). & \rk{2} & \rk{1}\tabularnewline[2.6pt]
\rr \textbf{Kick and well times} $1/\gamma_k$, $1/\gamma_s$ (Eq.~\ref{eq:pic-langevin}) & \rr $\approx7$ and $\approx100$ calls; ratio 14.8 [7.2, 30.4] (90\%), \#51 only (H130) & \rr a message acts for about 7 calls; an agent drifts back over about 100 & \rr To make an instruction hold, repeat it within 7 calls.\nt & \rk{1} & \rk{2}\tabularnewline[2.6pt]
\grp{Collective structure: do groups move as one?}
\rr \textbf{Collective modes}: $\lambda_1$ over the noise edge; participation ratio PR (Eq.~\ref{eq:spike}) & \rr above the edge: talk 21/24, content 24/24, activity 7/24 (H12, fragile); PR $-25\%$ in debates & \rr an activity mode above the trimmed edge: a shared drive besides the schedule & \rr Find the drive. A falling PR can mean a restatement loop (one point posted again and again). & \rk{2} & \rk{1}\tabularnewline[2.6pt]
\rr \textbf{Room contrast} $Q$: content split of rooms beyond random groups & \rr $Q=1.8$--5.7, significant in 5/7 periods (H100) & \rr rooms differ: each room is its own domain & \rr Pass information across rooms on purpose.\nt & \rk{2} & \rk{1}\tabularnewline[2.6pt]
\rr \textbf{Mode lifetimes} $\tau_k$: how long slow patterns last & \rr not measured; one slow mode, $\tau\approx23$--28\,d, regime I (H81) & \rr shows what persists, not only what moves together & \rr No action yet. & \rk{1} & \rk{1}\tabularnewline[2.6pt]
\bottomrule
\end{tabular}
\end{table*}

\begin{figure}[t]
\centering\includegraphics[width=\columnwidth]{js/s2_readout_example.pdf}
\caption{A message acts at the read-out call of its reader (model 02, H08). One G38 message (star) and the calls of its room-mates (bars). Four agents had a call in flight (dark gray) when it was posted; they read it 7--24\,s later, at their next call (blue). Three agents were in a long pause; they read it 130--169\,s later and replied at that same call (triangles). Figure~\ref{fig:readoutjump} shows the same step in 17 periods.}
\label{fig:readout}
\end{figure}

% ------------------------------------------------------------------ 02 + 09
\subsection{Agents are kinetic spins that couple at the read-out call (models 02 and 09)}\label{sec:p02}
\begin{figure}[t]
\centering\includegraphics[width=\columnwidth]{js/s2_kickoff_matrix.pdf}
\caption{The kickoff text is a field vector (model 11, H54). Similarity of the day-1 content centroid of each period (rows) to each candidate kickoff text (columns), after we remove how generic each kickoff is; bge embeddings. Orange: more alike than average; purple: less. Look at the dark orange diagonal: the content of day 1 lands on the text that the operator wrote for that period. The own kickoff ranks first in 18/33 periods (Fig.~\ref{fig:kickoffrank}).}
\label{fig:kickoff}
\end{figure}

At each call, an agent talks (it posts a message), works or waits. Let $Y_{ic}=1$ when call $c$ of agent $i$ is a talk call. A \emph{kinetic Ising model} describes a magnet in which each spin flips now and then, with a chance set by the field and by its neighbors. In our read-out version, an agent can flip only at its own calls, and only the messages that the call reads count as neighbors~\cite{glauber1963,aguilera2026}:
\begin{equation}
P(Y_{ic}=1)=\sigma\Big(\beta\Big[h_i(t_{ic})+\!\!\sum_{m\in B_{ic}}\!\!\big(J_0+J_{\rm N}\,\mathbf 1_{m\to i}\big)\Big]\Big).
\label{eq:pic-kinetic}
\end{equation}

\emph{Meaning.} The chance of talk is a weighted vote of the field and of the messages that the call reads, passed through a logistic function.
\begin{tightitem}
\item $h_i(t_{ic})$: the \emph{field} on agent $i$ at the call (scheduler, goal, room, habit);
\item $B_{ic}$: the messages that the call reads for the first time;
\item $J_0$: the \emph{coupling} per message read; $J_{\rm N}$: the extra coupling when $m$ names $i$ ($\mathbf 1_{m\to i}=1$);
\item $\beta$: the inverse temperature, how strictly the agent follows the vote (heat jostles the spins of a magnet; a hot agent, small $\beta$, acts nearly at random);
\item $\sigma(u)=1/(1+e^{-u})$: it changes the vote to a probability.
\end{tightitem}
\emph{Picture.} A person in a meeting decides at each turn whether to speak. The schedule and the habits of the person set a base level. Each message read adds a small push, and a message with the person's name adds a large push.

\emph{From the logs.} The \emph{jump per read} is the chance of talk at the read-out call minus the chance at the call in flight, at the same time since posting. The data fix only the products $\beta h$ and $\beta J$, not $\beta$ alone, so we report jumps.

\emph{Evidence.} A read that names the reader adds 0.079 [0.069, 0.089] talk calls; an unnamed read adds 0.004 [0.003, 0.006], about 19 times less (H67). Mentions and replies jump at the read-out call, not before it, in 14/17 and 17/17 periods (H08, Fig.~\ref{fig:readout}). Content also moves, five times more for named messages (H50).

\textbf{The clock is the call.} On the \emph{call clock}, time is the count of an agent's own calls: an agent that waits ten minutes between calls has not aged. To test it, note that when a call takes twice as long, the chance of a reply changes by a factor $2^{\eta}$. Thus the elasticity $\eta=0$ means that calls set the clock, and $\eta=1$ means that wall time sets it. In regimes II and III, $\eta=-0.11$ [$-0.16$, $-0.06$] (H40); regime I gives $+0.51$ [0.41, 0.61], which we cannot explain.

\textbf{Loop gain and its echoes.} Model 09 is a \emph{Hawkes process}: a model of events that trigger later events, first used for earthquake aftershocks~\cite{hawkes1971}. It counts the echoes of one message:
\begin{equation}
\begin{gathered}
g=\bar m\,\bar r\,J_1^\ast,\qquad \chi=\frac{1}{1-g},\\
\Phi=\frac{\mathrm{Var}\big(\sum_i n_i\big)}{\sum_i\mathrm{Var}(n_i)}=\frac{1}{(1-g)^2}.
\end{gathered}
\label{eq:fano}
\end{equation}

\emph{Meaning of $g$.} The \emph{loop gain} $g$ is the number of extra messages that one message causes, on average, at the first step. One message reaches $\bar r$ readers; each reader talks with an extra chance $J_1^\ast$ (the jump over the placebo); each talk call posts $\bar m$ messages.

\emph{Meaning of $\chi$.} The \emph{amplification} $\chi$ is the total response to one kick, with all echoes: $1+g+g^2+\dots=1/(1-g)$. It goes to infinity as $g\to1$, as the susceptibility of a magnet (its response to a small field) does when its temperature $T$ goes to the critical temperature $T_c$. Thus $1-g$ plays the role of the distance $T/T_c-1$ to the critical point.

\emph{Meaning of $\Phi$.} The \emph{Fano factor} of a count is its variance over its mean. It is 1 for events that arrive purely at random (a Poisson stream) and larger for bursty events. Our \emph{collective Fano ratio} $\Phi$ divides the variance of the group total of talk counts $n_i$ by the sum of the agent variances. It is 1 for independent agents. The echo loop multiplies each swing of the group total by $\chi$, and a variance grows as the square of a swing. Echoes add little to the variance of one agent, because they spread over many agents. Thus $\Phi=\chi^2$, with no free parameter.

\emph{Picture.} A microphone near a speaker gives a short echo at low volume ($g<1$: echoes shrink) and a howl at high volume ($g\ge1$).

\emph{From the logs.} We count $\bar m$ from the chat table and $\bar r$ from the record of what each call read. For $\Phi$ we count talk per agent in 10-min windows, without the day edges.

\emph{Evidence.} The median gain is $g=0.13$ (pooled 0.127 over 25 units); the upper bound is below 0.40 in 66/66 units (H67). Thus $\chi\approx1.15$ and $\Phi_{\rm pred}\approx1.32$ (both derived), and $\chi\le1.31$ in 33/33 periods (H51). The observed $\Phi$ matches: $r_F=\Phi_{\rm obs}/\Phi_{\rm pred}=1.01$ [0.94, 1.09] over 18 units (H111, Fig.~\ref{fig:fano}), on the same data as $g$. Regime I has no measured gain, and talk varies 1.34 [1.18, 1.52] times more than the loop permits.

\textbf{Variant: the context holds a self-field.} The recent output of an agent stays in its context and acts as a field on the agent. Thus \emph{restatement loops}, in which an agent posts the same point again and again, copy what the context holds (H69). In \#51 (each agent had a private goal), a forced erasure inside an \emph{idle trap}, a long run of waits, lifts the chance of a lasting escape from 0.47 to 0.84 (H16). Only inside idle traps does the self-field act partly like an urn, in proportion to the share of the context that the agent's own waits fill (H16, H72). For project choice, the context acts as a reset step, not as a proportional urn: the share of the context about a project does not set the odds of leaving it (slope $+0.07$ [$-0.05$, $+0.19$]; an urn predicts 1), but a forced reset raises them ($+0.40$ [$+0.25$, $+0.54$] log-odds, G51; H134, Sec.~\ref{sec:m02}).

% ------------------------------------------------------------------ 11 + 01
\subsection{External fields set most of the state (models 11 and 01)}\label{sec:p11}
\begin{figure}[t]
\includegraphics[width=\columnwidth]{js/m16_relaxation.pdf}
\caption{A fast read kick on a slow, overdamped well (model 16). (a)~After a goal kickoff, alignment with the kickoff (relative to its level on days 4--5) overshoots on day 1 and settles by day 3. Nothing dips in the gray undershoot window ($U=-0.006$ [$-0.020$, $+0.008$], bge, 27 kickoffs, 90\% CI; H125). Gray squares: the same with gte embeddings. (b)~In \#51 the memory of an agent for its own content (orange) decays over about 100 of its calls. The pull of a read (blue) decays in about 7 (H130). We divide each series by its fitted amplitude; thin lines are $e^{-\gamma n}$. (c)~Simulation: state $x$ = slow well $s$ + short kicks at reads (ticks). Look at the gap between the two curves in (b): one rate cannot describe both.}
\label{fig:m16}
\end{figure}

\begin{figure}[t]
\includegraphics[width=\columnwidth]{js/m17_modes.pdf}
\caption{Collective modes against a calibrated random-matrix edge (model 17, H12). (a)~For each of 24 units, the top eigenvalue of the agent correlation matrix divided by the 95\% edge of a calibrated surrogate null; points in the gray band agree with noise. Activity mostly fails the trimmed block-shift null (7/24). Talk (21/24) and content (24/24, median ratio 1.52) carry a real mode. (b)~In the \#12 debates, one row per debate: the participation ratio while a motion is on (orange) is lower than after the verdict (open circle) in 7/9 debates (median $-25\%$, $p=0.010$). Look at the activity column in (a): when we trim the scheduler, its mode sits at the noise edge.}
\label{fig:m17}
\end{figure}

Model 11 treats the content of each agent as an arrow, a \emph{vector spin} $\mathbf s_i$. The arrow is the mean \emph{embedding} of the statements of the agent in a window, minus the regime mean and its style. (An embedding is a vector from a text model; similar meanings give near vectors.) As in a magnet, each arrangement of arrows has an energy $\mathcal H$, and arrangements of low energy are more probable, with weight $e^{-\beta\mathcal H}$:
\begin{equation}
\begin{gathered}
-\beta\mathcal H=\sum_{i<j}J_{ij}\,\mathbf s_i\cdot\mathbf s_j+\sum_i\mathbf h_i\cdot\mathbf s_i ,\\
\mathbf h_i=\mathbf h_{\rm sched}+h_g\hat{\mathbf g}+\mathbf h_{r(i)}+\mathbf h_i^{0}.
\end{gathered}
\label{eq:vector}
\end{equation}

\emph{Meaning.} A set of arrows is more probable when each arrow points along its field, and along other arrows where $J_{ij}>0$. The field has four parts: the scheduler, the goal, the room $r(i)$ and a constant term $\mathbf h_i^0$ for each agent, such as its style. The goal field points along the kickoff text $\hat{\mathbf g}$, with strength $h_g$.

\emph{Picture.} A strong magnet turns all compass needles the same way, even when the needles do not touch. When you see aligned needles, find the magnet first. Model 01, the binary inverse Ising model, infers fields and couplings from on/off activity alone; it gives the gain and the field share in a simpler form~\cite{schneidman2006,meshulam2019}.

\emph{From the logs.} The \emph{field share} $f$ is the part of the excess co-activation (agents active at the same time, beyond chance) that the daily start and stop of the scheduler explain. We test it against a null that shifts the record of each agent in time by whole blocks: this keeps each agent's habits but breaks the timing between agents. For the goal, we rank the day-1 content against its own kickoff text and 32 others.

\emph{Evidence.} The daily start and stop give 70--80\% of the co-activation ($f=0.72$ median over 8 periods; regime I 0.11; H38, fragile). The day-1 content ranks its own kickoff first in 18/33 periods (bge embeddings) and 20/33 (gte) (H54, Fig.~\ref{fig:kickoff}; bge and gte are two text-embedding models). Style identifies an agent at 0.85, against a chance level of 0.15 (H46), and rooms split content in 5/7 periods (H100). A goal switch is a \emph{quench}, a sudden change of the field. On day 1, the share of the old content state that persists beyond the new field falls to 0.18--0.27, against 0.85--0.89 across a night (H96).

% ------------------------------------------------------------------ 16
\subsection{Each agent relaxes as a fast kick on a slow well (model 16)}\label{sec:p16}

Let the content of agent $i$ at its call $n$ be $x_i=s_i+k_i$, a slow part plus a fast part. A \emph{Langevin model} describes a particle that random kicks push and friction pulls back. Our two-rate version~\cite{uhlenbeck1930} gives
\begin{align}
s_i(n+1)-c_i&=(1-\gamma_s)\,\big(s_i(n)-c_i\big)+\xi_i(n),\nonumber\\
k_i(n+1)&=(1-\gamma_k)\,k_i(n)+\kappa\!\!\sum_{m\in B_{in}}\!\!(z_m-x_i)_\perp .
\label{eq:pic-langevin}
\end{align}

\emph{Meaning.} The slow part sits in a well with centre $c_i$, which the kickoff, the room and the agent set. At each call it moves back a fraction $\gamma_s$ of the way to the centre, and the noise $\xi_i$ pushes it at random. Each message $m$ that the call reads pulls the fast part toward the message content $z_m$, with strength $\kappa$. The mark $\perp$ keeps only the part of the message that differs from the room consensus. The fast part fades by a fraction $\gamma_k$ at each call. In $1/\gamma$ calls, a mark fades to about one third of its size.

\emph{Picture.} Recall the ball in a bowl of honey (Sec.~\ref{sec:intro}): the tap fades in a few calls, and the return to the bottom takes many more. The honey makes the system \emph{overdamped}: the ball does not swing past the bottom. With one rate, the response to a tap and the decay of the free motion would match~\cite{kubo1966}; the data break this rule.

\emph{From the logs.} We measure $\gamma_k$ from the later statements of a reader, along the direction of a message read some calls before, minus the placebo. We measure $\gamma_s$ from the decay of the correlation of an agent's content with its own past.

\emph{Evidence.} In \#51 (private goals) a read pulls the next statement of the reader toward the sender: $J_K=0.044$ [0.038, 0.050], in 11/12 units (H130). The pull fades in about 7 calls ($\gamma_k\approx0.13$--0.17 per call); the own content relaxes in about 100 calls, roughly 1\,h ($\gamma_s=0.0094$ [0.0076, 0.0115]). The ratio is 14.8 [7.2, 30.4] (90\% CI). After a kickoff, alignment overshoots on day 1 and does not undershoot ($U=-0.006$ [$-0.020$, 0.008], 90\% CI; H125; Fig.~\ref{fig:m16}). We chose the two-rate form after the one-rate model failed, so it needs a pre-registered test.

% ------------------------------------------------------------------ 17
\subsection{Beyond one shared mode, few collective patterns rise above the noise (model 17)}\label{sec:p17}

Take $N$ agents and $T$ time samples in a window. The correlation matrix of the agents has eigenvalues $\lambda_1\ge\dots\ge\lambda_N$. Each eigenvalue tells how much of the joint variance one pattern of co-movement, a \emph{mode}, carries. Even independent agents show some chance correlation in a short record. Random-matrix theory gives this noise floor, the \emph{Marchenko--Pastur edge} $\lambda_+$, and the condition for a real mode~\cite{marchenko1967,baik2005}:
\begin{equation}
\lambda_+=\big(1+\sqrt q\,\big)^2,\quad q=\frac NT;\qquad
\text{mode real if}\ \ell>\sqrt q .
\label{eq:spike}
\end{equation}

\emph{Meaning.} For independent agents, all eigenvalues stay below the edge $\lambda_+$. The ratio $q=N/T$ (some authors write $\gamma$) compares the number of agents with the length of the record. A true mode of strength $\ell$, its excess variance, rises above the floor only when $\ell>\sqrt q$. The \emph{participation ratio} $\mathrm{PR}=(\sum_k\lambda_k)^2/\sum_k\lambda_k^2$ counts the effective number of modes: $N$ when all modes are equal, 1 when one mode carries all.

\emph{Picture.} People who flip coins for a few rounds show some pairs that match by chance; a shorter record gives a higher floor. A real mode is a choir that must sing louder than the crowd.

\emph{From the logs.} We compute $\lambda_1$ from the agent-by-time matrix of activity, talk or content. Real series are autocorrelated, so we replace $\lambda_+$ by the 95th percentile of $\lambda_1$ for surrogates, copies of the data in which each agent is shifted in time~\cite{laloux1999}.

\emph{Evidence.} Talk carries a mode above this edge in 21/24 units and content in 24/24. Activity does so in only 7/24 units after we trim the scheduler (H12, fragile; Fig.~\ref{fig:m17}). In the \#12 debates (two teams of agents argued motions, and one agent judged), a motion lowers the PR of content by 25\% (7/9 debates). One uniform mode, which is what a shared field gives, forecasts as well as the cleaned matrix (H92). We have not yet measured the threshold $\ell>\sqrt q$ itself.

% ------------------------------------------------------------------ rest of picture + observables
\subsection{Memory, groups and the observables}\label{sec:observables}

Two results from partly useful models complete the picture (Sec.~\ref{sec:partial}). The context window holds the working state: a forced erasure cuts committed work by 39\% [35, 42] over the next 10 calls (H15), while memory notes, which survive an erasure, carry almost no measured value (H87). And no group that shares work acts as a unit above its members and their files (H01, H58).

The picture is also a design rule. On the evidence so far, the way to make a swarm coordinate is through the field: write the kickoff, name the agents you want to move, and keep their context windows intact. The chat between agents adds little at these sizes, and the chat-based levers (nudges, corrections) move almost nothing. If independence is the goal, the shared inputs are what to remove, not the chat. We have not tested these design rules as interventions; they are what the fitted picture implies.

Table~\ref{tab:monitor} turns the picture into 17 ranked observables, each one number per goal period or week. Appendix~\ref{app:observables} gives the full definitions and proposed alarm levels.

